% Compile twice with XeLaTeX. Slide numbers match review.tex.
\documentclass[a4paper,11pt]{article}
\usepackage{fontspec}
\usepackage[russian]{babel}
\setmainfont{Arial}
\usepackage[margin=20mm]{geometry}
\usepackage{amsmath,amssymb,hyperref}
\setlength{\parindent}{0pt}
\setlength{\parskip}{5.5pt}
\setlength{\emergencystretch}{3em}
\linespread{1.04}
\newcommand{\slide}[2]{\clearpage\section*{Слайд #1. #2}}
\newcommand{\question}[1]{\par\textbf{#1}\par}
\hypersetup{pdftitle={Подробная речь: VQ4DiT и DPQ},pdfauthor={},
  pdfsubject={Пояснения к 28 слайдам: устройство DiT, квантизация, две статьи и ответы на вопросы},
  pdfcreator={XeLaTeX},hidelinks}
\begin{document}

\slide{1}{Векторная квантизация DiT: VQ4DiT и DPQ}
Задача обеих статей --- уменьшить объём хранения диффузионного трансформера,
не разрушив его способность генерировать изображения. Исходная сеть уже
обучена. Поэтому речь идёт не о создании генератора с нуля, а о компактном
представлении его весов и последующей настройке этого представления.

Векторная квантизация заменяет несколько чисел одним коротким вектором
из небольшой таблицы --- кодбука. Вместо исходных чисел сохраняется номер
строки таблицы. Разные участки весов могут использовать одну строку, поэтому
таблица многократно переиспользуется. Возникают две задачи: подобрать значения
в таблице и определить, какой участок на какую строку ссылается.

VQ4DiT главным образом меняет вторую часть. На время обучения каждому
участку разрешается смешивать несколько кандидатов, а затем остаётся один.
DPQ сначала даёт разным группам столбцов независимые кодбуки, после чего
сжимает сами таблицы и чередует изменение индексов с обучением их значений.

Перед этими двумя методами необходимо пояснить, как устроен DiT и почему
его ошибка не ограничивается одним проходом. Одна и та же сеть вызывается
много раз при получении изображения. Поэтому хорошее приближение весов
ещё не гарантирует хорошую генерацию.

\question{Какие числа приведены дальше?}
Таблицы VQ4DiT и DPQ содержат опубликованные результаты соответствующих
статей, а не результаты кластерных репродукций проекта. Их назначение ---
показать возможности авторских методов и условия, при которых получены числа.
Для утверждения о воспроизведении нужны отдельные завершённые измерения.

\slide{2}{Один шаг генерации проходит через всю DiT}
DiT работает не непосредственно с RGB-пикселями. Изображение размером
$256\times256$ представляется латентом размером $32\times32$ с четырьмя
каналами. При генерации готовой картинки ещё нет: начальный латент сразу
заполняется случайным шумом. В конце отдельная сеть VAE переводит очищенный
латент в изображение.

Трансформеру нужна последовательность векторов. Поэтому латент делится
на непересекающиеся участки $2\times2$. В каждом находятся четыре позиции
по четыре канала, то есть 16 чисел. Отсюда получается $16\times16=256$
участков. Обучаемая проекция переводит каждый участок в 1152 признака,
добавляется информация о положении. Такие векторы называются токенами.

Все 256 токенов совместно проходят через 28 последовательных блоков DiT.
Выходная часть предсказывает шум и параметры дисперсии перехода.
Правило DDPM по этим ответам вычисляет следующий латент, затем этот
латент снова проходит через все 28 блоков. В режиме DDPM50 выполняется
50 таких переходов. Это не 50 блоков и не 50 обучающих обновлений.

\question{Что означает «/2» в названии модели?}
Это размер участка в латентном пространстве. Не два пикселя исходного
изображения и не два диффузионных перехода. Между участком входных данных
и вектором весов, который будет квантуемым объектом, также нет тождества.

\question{Каждый токен отдельно вызывает сеть?}
Нет. Все токены обрабатываются одним вызовом. Attention позволяет им
обмениваться информацией. При следующем времени вызывается та же самая
сеть с теми же весами, но уже с новым латентом и другим значением времени.
Исходная сетка содержит 1000 времён; DDPM50 использует сокращённую сетку
с пересчитанными коэффициентами переходов.

\slide{3}{Слои одного блока и управление условием}
В каждом блоке DiT находятся два основных преобразования: attention и MLP.
Attention связывает разные участки изображения. Один участок получает
информацию о других, поэтому сеть может согласовывать положение и части
объекта. MLP отдельно преобразует признаки каждого токена, используя
расширение размерности, нелинейность и обратное сжатие.

У attention рассматриваются две матрицы. Первая одновременно получает
три набора признаков Q, K и V: $1152\to3456$, около 3.98 млн весов.
Вторая переводит результат обратно в 1152 признака и содержит около
1.33 млн весов. Две матрицы MLP имеют размеры $1152\to4608$ и
$4608\to1152$, по 5.31 млн весов каждая.

Пятая матрица --- adaLN, $1152\to6912$, около 7.96 млн весов.
Она получает условие, составленное из времени и класса, и выдаёт
шесть векторов по 1152 числа: сдвиг, масштаб и управляющий множитель
для attention, а также такие же три вектора для MLP. Поэтому adaLN
не является маленькой несущественной поправкой: она содержит треть
исследуемых весов блока и управляет обеими ветками.

Всего получается около 23.89 млн весов в пяти матрицах одного блока.
У 28 блоков свои параметры, поэтому область квантизации содержит
140 матриц и 668.86 млн весов. LayerNorm, нелинейности и остаточные
связи продолжают работать как обычно: они не превращаются в кодбуки.

\question{Почему на выходе восемь каналов?}
Четыре относятся к предсказанию шума латента, ещё четыре --- к обучаемой
дисперсии перехода. Именно выходная часть после всех блоков формирует
эти ответы. В рассматриваемой области сжатия она остаётся полноточной.

\question{Квантуются ли активации?}
Нет. Сжимаются выбранные матрицы весов. Числа, которые проходят через сеть
при обработке изображения, сохраняют представление с плавающей точкой.
Учёт активаций в ошибке и квантизация самих активаций --- разные действия.

\slide{4}{Что меняется при переходе к другому времени}
В генерации нужно различать глубину модели и время диффузии.
Глубина --- 28 последовательно расположенных блоков. Каждый имеет свои
матрицы attention, MLP и adaLN. Это не повторное использование одной
матрицы 28 раз. Время --- номер состояния, с которым вся сеть вызывается
в очередной раз. При смене времени веса не заменяются.

Число времени сначала переводится в синусоидальные признаки, затем
небольшая MLP получает из них обученный вектор. Номер класса выбирает
строку другой обученной таблицы. Эти два вектора складываются:
$c=e_t(t)+e_y(y)$. Полученное условие поступает в каждый блок.
Однако матрица adaLN у каждого блока собственная, поэтому одинаковое
условие преобразуется в разные управляющие значения.

Сдвиг прибавляет величину к нормированным признакам. Масштаб меняет
величину признаков. Управляющий множитель определяет вклад результата
ветки в остаточное обновление. При другом времени и другом латенте
получаются другие активации, хотя все параметры сети остаются прежними.

\question{Нужны ли разные кодбуки для разных времён?}
В VQ4DiT и DPQ этого нет. Кодбуки принадлежат матрицам и используются
при всех временах. Градиенты от разных времён помогают подобрать один
набор значений, который должен работать на всём нужном диапазоне.
Временная специализация таблиц была бы отдельным изменением алгоритма.

\question{Что остаётся вне области сжатия?}
Кодировщики времени и класса, входная проекция и выходной модуль.
Выходная часть содержит около 2.69 млн параметров. Её малый размер
относительно 668.86 млн исследуемых весов позволяет сохранить точность
этой части без основной потери экономии памяти.

\slide{5}{Как шум превращается в изображение}
Сначала создаётся случайный латент. Сеть ещё не видит осмысленного
объекта: на входе шум, номер времени и класс. После прохода DiT
правило перехода получает оценку шума и дисперсии и строит следующий
латент. Новый латент содержит немного другую информацию и становится
входом следующего вызова. Так постепенно появляется структура изображения.

Цепочка всех промежуточных состояний называется траекторией генерации.
Это просто последовательность массивов $z^{(0)},z^{(1)},\ldots,z^{(50)}$.
Латенты внутри цепочки не обязаны выглядеть как готовые картинки при
декодировании: на ранних этапах шум ещё очень велик.

Если заменить веса сжатыми, ответ одного вызова немного изменится.
Изменится следующий латент, а значит, на следующем переходе сеть
получит уже другой вход. Поэтому ошибка может не только добавляться,
но и усиливаться или частично компенсироваться. Сравнение двух сетей
на одном входе и сравнение двух полных траекторий --- разные опыты.

\question{Откуда VQ4DiT получает обучающие входы?}
Исходная сеть выполняет эту цепочку. Из неё берутся промежуточные
латенты вместе с соответствующими временами и классами. По ним можно
сравнить исходную и сжатую DiT. Для этого не требуется набор реальных
изображений: источником примеров выступает готовый генератор.

\question{Итоговое изображение становится следующей обучающей целью?}
Не в этом методе. Целью VQ4DiT служат выходы внутренних блоков исходной
сети на выбранных промежуточных состояниях. Траектория даёт входы,
а проход исходной DiT по ним даёт эталонные признаки.
После калибровки сжатая модель должна генерировать самостоятельно,
без помощи исходной сети в каждом переходе.

\slide{6}{ODE и SDE: два способа обновлять латент}
Нейросеть и правило генерации выполняют разные задачи. Нейросеть
предсказывает, как интерпретировать текущий шумный латент.
Семплер --- алгоритм, который использует ответ сети и вычисляет
следующий латент. ODE и SDE относятся к математическому описанию
этого движения, а не к отдельному слою трансформера.

SDE --- стохастическое дифференциальное уравнение: во время движения
появляются новые случайные добавки. В используемой DDPM-генерации
это видно в переходе как член $\sigma_t\xi_t$, где $\xi_t$ --- новый
гауссовский шум. На последнем переходе добавка отключается.
Даже при одинаковом начальном латенте другие случайные добавки
могут привести к другой траектории и другой картинке.

ODE --- обыкновенное дифференциальное уравнение. После задания
начального состояния его движение не требует новых случайных толчков.
Связанный практический пример --- детерминированный DDIM с нулевым eta.
Случайным остаётся начальный латент, но затем фиксированные модель,
класс и правило перехода задают путь.

\question{Можно просто занулить шум в DDPM и назвать это ODE?}
Нет. DDPM и DDIM являются дискретными правилами перехода, а ODE/SDE ---
непрерывными описаниями. Корректный детерминированный переход имеет
другие коэффициенты и преобразования. Простое удаление случайного
слагаемого не воспроизводит автоматически DDIM или probability-flow ODE.

\question{SDE используется при калибровке, а ODE при оценке?}
Такого разделения нет. В рассматриваемом DDPM-протоколе стохастическая
цепочка используется и при создании входов исходной моделью, и при
итоговой генерации сжатой моделью. Калибровка и оценка --- цели запуска,
а не названия разных типов движения.

Отсутствие нового шума не гарантирует более высокий FID или меньший FID,
меньшую ошибку квантизации либо ускорение. Основное время тратится
на вызовы DiT. Для суждения о качестве требуется отдельное сравнение
при совпадающем числе переходов и остальных условиях.

\slide{7}{Калибровка и оценка: что запускается}
Калибровка --- настройка компактного представления готовой сети.
Оценка --- независимая проверка, какие изображения получаются после
этой настройки. Для обоих методов в оценке запускается уже сжатая DiT,
которая сама проходит всю цепочку генерации.

В VQ4DiT обучающие входы получаются из последовательной генерации
исходной DiT. Например, после нескольких переходов существует
латент $z_t$. Он подаётся исходной и сжатой сетям, затем сравниваются
их внутренние признаки. Следующее состояние обучающей цепочки
продолжает задавать исходная модель. Это защищает источник обучающих
входов от накапливающихся ошибок ещё не настроенной сжатой сети.

В DPQ есть реальные изображения ImageNet. VAE переводит изображение
в чистый латент. Затем выбирается одно время и сразу добавляется
известный шум нужной величины. Получается один обучающий пример.
Не требуется сначала выполнять предыдущие обратные переходы,
чтобы попасть на выбранное время. Цель при обучении кодбуков
связана с известным добавленным шумом.

\question{Значит ли прямое зашумление, что DPQ обучается на SDE-траектории?}
Нет. Случайный массив создаётся сразу для одного времени. Это выборка
из распределения прямого зашумления, а не численное прохождение
стохастической цепочки обратной генерации.

\question{Что делает текущий перенос VQDM?}
Исходная DiT в режиме DDPM50 даёт входы линейных слоёв. Их статистика
используется при локальном подборе кодбуков и индексов.
При оценке квантованная сеть сама выполняет DDPM50.
Это отличается и от сравнения всех промежуточных выходов VQ4DiT,
и от задачи предсказания известного шума в M-шаге DPQ.

\slide{8}{На что влияет выбор способа генерации}
Смена семплера не меняет архитектуру DiT и не создаёт новые кодбуки.
Однако она меняет последовательность входных латентов, поэтому
способна заметно влиять на изображения и на поведение сжатой модели.

Первое следствие --- метрики качества. Даже одна и та же исходная
сеть может иметь разные FID при DDPM50 и DDPM250. Поэтому результат
сжатой сети на 50 переходах нельзя непосредственно противопоставить
другому методу на 250 и приписать всё различие квантизации.

Второе следствие касается калибровочных входов. Если они собираются
из генерации, то другой семплер даёт другие состояния и активации.
Неизменные кодбуки могут работать на них иначе. В DPQ обучающие
времена выбираются при прямом зашумлении, но качество на конкретном
обратном генераторе всё равно требует проверки.

Третье --- повторяемость. Для двух DDPM-генераций нужно согласовать
класс, начальный латент и случайные массивы всех переходов.
Одного одинакового начального массива недостаточно, если последующие
шумы различаются. У детерминированного семплера достаточно общего
старта при зафиксированных модели, условии и численных настройках.

\question{Детерминированный путь быстрее?}
Не обязательно. Генерация случайного массива обычно намного дешевле
прохода большой DiT. Число вызовов сети определяет основную стоимость.
Ускорение нужно измерять, а не выводить из отсутствия случайной добавки.

\question{Что фиксировать для честной таблицы?}
Модель, число переходов, правило обновления, CFG, VAE, число изображений,
эталонные статистики и оценщик. Отдельно нужно учитывать фактический
объём хранения, а не только красивое название битного профиля.

\slide{9}{VQ4DiT: зачем обучать назначения}
Первая статья --- \href{https://ojs.aaai.org/index.php/AAAI/article/view/33782/35937}
{VQ4DiT: Efficient Post-Training Vector Quantization for Diffusion Transformers},
AAAI 2025. Её основной вопрос: достаточно ли менять значения кодбука,
если закреплённые назначения весов могут быть неудачными?

При обычной VQ несколько групп весов ссылаются на один кодовый вектор.
Во время обучения каждая группа посылает градиент в одну и ту же строку
кодбука. Эти градиенты складываются. Если разные группы требуют
противоположных исправлений, итоговое изменение оказывается компромиссом.

Простой условный пример: один участок хочет сдвинуть общий центр вправо
на 0.1, другой --- влево на 0.1. При одинаковой силе сигналов их сумма
равна нулю. Перемещение единственного общего центра не может независимо
исправить оба участка. Для одного из них полезнее выбрать другую строку
кодбука, а не продолжать менять общую.

VQ4DiT сохраняет небольшое число кандидатов для каждого участка весов.
Во время калибровки допускается непрерывная смесь этих кандидатов.
Сеть учится одновременно менять значения в кодбуке и предпочтение
кандидатов. Затем остаётся один индекс, чтобы не увеличивать постоянное
представление весов двумя долями и двумя номерами.

\question{Почему не выбрать центр один раз по расстоянию весов?}
Близость весов --- полезная инициализация, но не вся задача.
Две похожие ошибки весов могут по-разному влиять на реальные активации
и выходы блоков. Поэтому окончательное предпочтение определяется
не только евклидовым расстоянием, а обучающей целью по признакам исходной сети.

\question{Все назначения пересматриваются среди всех центров?}
Нет. В основном варианте для каждого участка заранее выбраны два
кандидата. Позднейшее обучение выбирает между ними. Это важное
ограничение выразительности и вычислительной стоимости метода.

\slide{10}{Что хранится вместо матрицы весов}
Линейный слой хранит матрицу $W$ с $O$ строками и $I$ столбцами.
Каждая строка содержит веса, которыми входные признаки объединяются
для получения одного выходного признака. VQ4DiT делит строку
на соседние фрагменты длины $d$. Каждый фрагмент заменяется одной
строкой общего для этой матрицы кодбука.

Если $K=256$, номер строки таблицы помещается в восемь бит.
При длине фрагмента четыре числа один восьмибитный номер кодирует
четыре веса: номинально два бита на вес. Если $K=64$ и $d=2$,
шесть бит номера делятся на два веса: получается три бита на вес.
Формула проста: $b=\log_2(K)/d$.

Для матрицы $1152\times1152$ при $d=4$ нужно 331\,776 индексов.
Кодбук при этом всего один: 256 строк по четыре числа. Один и тот же
вектор таблицы могут использовать фрагменты разных строк и разных
позиций столбцов этой матрицы. Соседняя матрица имеет отдельную таблицу.

В пяти матрицах каждого из 28 блоков получается 140 независимых кодбуков.
Это отличается от аддитивного VQDM с четырьмя таблицами на матрицу:
в VQ4DiT нет суммы четырёх кодовых векторов для одного участка.
После калибровки используется один выбранный вектор.

\question{Реальный размер точно равен двум битам на вес?}
Нет. Два бита относятся только к индексам. Дополнительно хранятся
кодбуки, смещения и несжатые модули сети. Для честного сравнения
нужен общий размер модели либо битность с учётом всех этих частей.

\question{Набор таблиц дублируется для каждого времени?}
Нет. У матрицы один кодбук для всех времен. Когда меняются время,
класс и латент, меняются данные вычисления, а не список кодовых векторов.

\slide{11}{Инициализация: k-means и два кандидата}
До обучения нужно получить разумное компактное приближение весов.
Фрагменты одной матрицы рассматриваются как точки в пространстве
размерности $d$. Например, при $d=4$ каждый фрагмент --- четвёрка чисел.
Алгоритм k-means группирует близкие фрагменты и вычисляет центры групп.
Эти центры становятся начальными кодовыми векторами.

Смысл k-means в данном случае --- уменьшить ошибку приближения исходных
весов. Это ещё не обучение по генерации и не сравнение изображений.
После получения таблицы для каждого фрагмента находятся два ближайших
центра по евклидовому расстоянию. Их номера фиксируются как кандидаты.
Начальное предпочтение одинаковое: половина первого и половина второго.

Условный двумерный пример помогает отделить номера от значений.
Пусть целевой фрагмент равен $(1,1)$, а кандидаты --- строки с номерами
7 и 19. Их численные значения могут быть $(0.8,1.0)$ и $(1.1,1.1)$.
Начальная смесь равна $(0.95,1.05)$. Номера 7 и 19 остаются кандидатами,
но их значения и доли могут измениться при дальнейшей калибровке.
Это пример механизма, не числовой результат статьи.

\question{Почему два кандидата, а не все 256?}
Хранить и обучать 256 долей для каждого из сотен миллионов фрагментов
слишком дорого. Два кандидата оставляют возможность изменить назначение,
но сохраняют небольшой дополнительный объём на время обучения.
Число кандидатов далее проверяется в отдельной таблице статьи.

\question{Если один кандидат стал плохим после перемещения центров?}
Полный список всех центров заново не открывается в каждой итерации.
Это ограничение метода. Он рассчитывает, что близкие по весам кандидаты
дадут достаточно полезных вариантов, а обучение по работе сети
уточнит выбор внутри этой небольшой области.

\slide{12}{VQ4DiT: калибровка по выходам блоков}
Сначала исходная DiT запускает генерацию из случайного латента.
На промежуточном переходе появляется состояние $z_t$.
Оно вместе со временем и классом становится общим входом для
исходной и квантованной сетей. Реальная картинка ImageNet
для такого входа не требуется: источник данных --- готовый генератор.

В исходной сети сохраняются выходы всех 28 блоков.
Квантованная сеть тоже проходит свои 28 блоков на том же исходном
латенте. Получаются 28 пар массивов признаков: после первого блока,
после второго и так далее. Разницы между соответствующими массивами
задают обучающий сигнал.

Исходная модель фиксирована. В сжатой меняются значения кодбуков
и доли кандидатов, чтобы её внутренние признаки стали ближе
к исходным. В рассматриваемой реализации в обучаемые группы также
входят смещения блоков; это не превращает цель в сравнение картинок.
Следующее состояние калибровочной генерации продолжает задавать
исходная сеть, а не ещё не настроенная сжатая.

\question{Входы всех внутренних блоков точно совпадают?}
Нет. Совпадает вход в полную сеть. После первого квантованного блока
его результат может отличаться от эталона, и следующий блок получает
уже изменённые признаки. Нельзя автоматически считать, что каждый
блок обучается на принудительно подставленном исходном внутреннем входе.

\question{Сколько выходов реально нужно хранить?}
На выбранном входе сравниваются выходы всех 28 блоков. Для батча 16
без удвоения CFG один FP32-выход формы $16\times256\times1152$
занимает около 18.9 МБ, а 28 --- около 528.5 МБ. Это расчёт одного
набора эталонных признаков, не всей обучающей памяти: градиенты,
промежуточные активации и дополнительные ветви требуют больше.
Метод не требует хранить все выходы всех возможных времён одновременно.

\slide{13}{Как градиент доходит до выбора индекса}
Целочисленный индекс не имеет обычной производной, позволяющей
оптимизатору плавно поменять 7 на 19. Поэтому на время обучения
используется непрерывное представление выбора: две доли кандидатов.
Они неотрицательны и в сумме равны единице.

Удобная реализация --- softmax двух обучаемых чисел, логитов.
При равных логитах доли равны $0.5/0.5$.
Восстановленный фрагмент весов --- взвешенная сумма двух кодовых
векторов. Это обычное дифференцируемое вычисление. Ошибка выхода
блока зависит от полученных весов, значит, градиент достигает
и значений кодбука, и логитов долей.

В условном примере доли могут стать $0.1/0.9$.
Тогда второй центр в девять раз сильнее влияет на текущий фрагмент,
хотя на начальной геометрической проверке он мог быть лишь вторым
по близости. Это не скачок целого индекса: изменилось непрерывное
предпочтение между двумя заранее выбранными номерами.

\question{Что именно фиксировано?}
Номера двух кандидатов. Координаты строк кодбука по этим номерам
могут обучаться. Доли тоже обучаются. После окончательного выбора
хранится один номер, а мягкая смесь больше не нужна.

\question{Одна итерация здесь равна одному переходу DDPM?}
Нет. Переход DDPM обновляет латент для генерации.
Итерация оптимизатора обновляет параметры компактной модели.
Они могут быть организованы вокруг одной калибровочной цепочки,
но это разные операции: первая меняет данные, вторая --- обучаемые числа.

В отличие от прямого округления с подставной производной,
в этой схеме обучаемая смесь действительно дифференцируема.
Основная трудность переносится на финальный переход от смеси
к одному центру: мягкое приближение должно стать достаточно уверенным.

\slide{14}{Цель: повторить выходы исходных блоков}
Цель VQ4DiT складывается из различий между выходами соответствующих
блоков. Обозначение $h_l^{\mathrm{FP}}$ означает массив признаков
после блока с номером $l$ в исходной сети, а $h_l^{\mathrm Q}$ ---
после соответствующего блока в квантованной. Сначала для каждого
из 28 блоков считается квадрат разницы, затем результаты объединяются.

Почему это содержательнее обычной ошибки весов? Выход блока
учитывает attention, обе матрицы MLP, нелинейности, нормализацию,
управление через adaLN и остаточные связи. Можно плохо приблизить
некоторые несущественные веса, но сохранить важные вычисления.
И наоборот: небольшая средняя ошибка весов может затронуть
чувствительное направление признаков и сильно испортить выход.

Кодбуки и доли меняются так, чтобы воспроизводить именно эти
внутренние представления. Близость к исходным весам определяет
хороший старт и набор кандидатов, но не заменяет обучающую цель.
При сравнении последовательных блоков учитывается и влияние
ошибок предыдущей части квантованной сети.

\question{Это дистилляция?}
Да, в смысле обучения по признакам фиксированного учителя.
Учительская модель не становится меньше; сжатая модель учится
повторять её внутренние ответы. Это не обучение напрямую
по реальным изображениям или по известному добавленному шуму.

\question{Это уже финальное качество генерации?}
Нет. Обучающий показатель относится к внутренним признакам
на выбранных входах. После калибровки требуется самостоятельная
генерация квантованной сетью и оценка FID и других метрик.
Даже точные промежуточные признаки не являются математической
гарантией совпадения всего распределения изображений.

\slide{15}{Как смесь превращается в один индекс}
Мягкая смесь удобна для градиентного обучения, но не подходит как
окончательное компактное представление: пришлось бы сохранять два
индекса и две доли для каждого фрагмента. Поэтому нужно сделать
предпочтение уверенным и в конце оставить одного кандидата.

Для этого к ошибке блоков добавляется штраф, который велик
при неопределённых долях и мал около нуля и единицы.
Для отдельной доли используется выражение $1-|2r-1|$.
При $r=0.5$ оно равно единице, при $r=0$ или $r=1$ --- нулю.
Штраф суммируется по кандидатам и усредняется по фрагментам весов.

Для двух кандидатов доли $0.5/0.5$ дают сумму два на фрагмент,
$0.1/0.9$ --- 0.4, а $0/1$ --- ноль. Это пояснение формулы,
а не новые измерения статьи. Ошибка признаков выбирает полезный
вариант, а дополнительный штраф подталкивает к дискретности.
Один штраф без сравнения с учителем не определил бы,
какой из двух центров лучше для генерации.

Когда штраф достаточно мал, в статье предлагается выбрать
кандидата с максимальной долей и перестать обучать доли.
Приведён пример порога $10^{-4}$. При итоговой генерации
фрагмент восстанавливается одной строкой кодбука, без смеси.

\question{Почему не взять максимум сразу?}
Жёсткий выбор с самого начала не позволяет обычному градиенту
плавно исследовать изменение назначения. Мягкий этап даёт
такую возможность. Если закончить его слишком рано, полезная
комбинация двух центров может исчезнуть при округлении выбора.

\question{Штраф сам разрывает абсолютно равные доли?}
Не обязательно. В симметричной точке модуль недифференцируем,
а типичная реализация выбирает нулевую производную в изломе.
Асимметрию обычно создаёт ошибка признаков, поскольку кандидаты
дают разные ответы. Поляризация усиливает возникшее предпочтение,
но не является самостоятельным поиском правильного индекса.

\slide{16}{Полный проход VQ4DiT}
Алгоритм начинается с готовых весов исходной DiT.
В каждой матрице они разбиваются на короткие фрагменты.
K-means получает начальный кодбук, а для каждого фрагмента
по расстоянию весов выбираются два кандидата.
Начальные доли одинаковы.

Далее исходная DiT создаёт последовательность промежуточных
латентов через DDPM. На выбранных состояниях исходная
и квантованная сети проходят свои блоки. Сравниваются
28 пар признаков. Ошибка признаков и штраф неопределённых
долей дают общий обучающий сигнал. Оптимизатор меняет
значения кодбуков и предпочтения кандидатов.

После калибровки остаётся один индекс на фрагмент.
Тогда уже квантованная модель самостоятельно генерирует
новые изображения. Исходная сеть и мягкие доли не должны
требоваться при таком использовании.

\question{Какие настройки приведены в статье?}
RMSprop, шаг $5\cdot10^{-2}$ для долей и $10^{-4}$ для
остальных обучаемых параметров, батч 16, 500 обновлений.
Это не 500 DDPM-переходов и не 500 эпох ImageNet.
В статье указана одна A100 и время от 20 минут до пяти часов
для разных настроек. Эти времена нельзя выдавать за
измеренную стоимость текущей кластерной репродукции.

\question{Что в реализации нужно фиксировать особенно аккуратно?}
Обучаются ли смещения блоков, как суммируются ошибки разных
блоков и штрафы матриц, когда выполняется жёсткий выбор,
какое число обновлений реально используется. В опубликованном
коде и текстовом описании есть различия; они записаны
отдельно в протоколе воспроизведения. Также статья не вполне
уточняет, включает ли батч 16 удвоение условной и безусловной
ветвей CFG. Поэтому каждая репродукция должна явно
зафиксировать свой выбор, а не скрывать его за названием метода.

\slide{17}{VQ4DiT: опубликованные результаты}
Все числа этой таблицы взяты из таблицы 3 статьи.
Модель --- DiT-XL/2, изображения $256\times256$, CFG 1.5,
для каждой оценки генерируются 10\,000 изображений.
Строки с 50 и 250 переходами нужно рассматривать отдельно.

FID измеряет различие статистик признаков реальных
и сгенерированных изображений: меньше лучше.
sFID использует пространственную информацию признаков
и тоже уменьшается при лучшем согласовании.
IS оценивает уверенность распознавания отдельных изображений
и разнообразие предсказываемых классов. Precision отражает,
насколько сгенерированные примеры попадают в область
реалистичных данных. Последние две метрики направлены вверх.

На 50 переходах исходная модель имеет FID 6.72.
У трёхбитного VQ4DiT FID 11.91, у двухбитного --- 12.42.
На 250 переходах исходная модель даёт 5.33,
трёхбитная --- 10.59, двухбитная --- 11.87.
То есть генерация не разрушается даже при двух битах
на индексы, но разрыв с исходной моделью остаётся.

\question{Почему IS или precision иногда выше исходной модели?}
Метрики оценивают разные свойства. Генератор может выдавать
более уверенно распознаваемые и типичные примеры,
но хуже покрывать распределение реальных изображений.
Поэтому высокий IS или precision не отменяет худший FID
и не доказывает, что сжатая модель в целом лучше исходной.

\question{Можно сравнить 11.87 со значением собственного VQDM?}
Только после согласования протокола: числа переходов,
VAE, оценщика, количества примеров, условного усиления
и фактической памяти. Номинальные два бита VQ4DiT
не являются тем же бюджетом, что около 4.15 бита
всей DiT в текущем VQDM. Сначала важно сравнить
каждый метод со своим корректным исходным контролем.

\slide{18}{Что даёт свобода выбора назначения}
Это отдельная проверка идеи VQ4DiT из таблицы 5 статьи.
Архитектура, двухбитный профиль и генерация DDPM50
сохраняются. Меняется число кандидатов, которые разрешены
для каждого фрагмента весов во время калибровки.

Один кандидат означает, что начальное назначение фиксировано.
Кодовые векторы ещё могут двигаться, но фрагмент не может
перейти на другую строку таблицы. В этом варианте FID равен 60.16.
Два кандидата дают FID 12.42, три --- 12.22,
четыре --- 14.06. Главный скачок наблюдается при появлении
самой возможности перевыбора, то есть от одного к двум.

\question{Это доказывает, что три кандидата всегда лучше двух?}
Нет. Показана небольшая разница в одном авторском сравнении.
Повторные обучения и интервалы неопределённости здесь не приведены.
Кроме того, больше кандидатов означает больше долей,
больше временной памяти и более сложное обучение выбора.
Дальнейшее увеличение не обязано давать монотонный выигрыш.

\question{Можно утверждать, что любой метод без перевыбора плох?}
Нет. Вариант с одним кандидатом имеет конкретную инициализацию,
обучающую цель и настройки этой статьи. Он показывает
важность назначения внутри данной процедуры,
но не заменяет отдельное сравнение с более сильной
подготовкой кодбуков, например с готовым VQDM.

Главный вывод для следующих исследований: значения таблицы
и назначения нельзя считать полностью независимыми решениями.
Даже хорошие кодовые векторы могут использоваться неудачно.
Однако VQ4DiT не открывает бесконечный поиск всех вариантов:
качество ограничено исходным небольшим набором кандидатов,
а мягкая оптимизация должна ещё пережить жёсткий финальный выбор.

\slide{19}{DPQ: от общего кодбука к подпространствам}
Вторая статья --- \href{https://openaccess.thecvf.com/content/ICCV2025/papers/Shao_Memory-Efficient_Generative_Models_via_Product_Quantization_ICCV_2025_paper.pdf}
{Memory-Efficient Generative Models via Product Quantization}, ICCV 2025.
Отправная гипотеза отличается от VQ4DiT: одной небольшой
таблицы для всех позиций матрицы может не хватать.
Разные входные признаки могут требовать разных наборов
типичных фрагментов весов.

В варианте с общим кодбуком первые четыре столбца,
следующие четыре и все остальные позиции выбирают
векторы из одного и того же набора. DPQ выделяет каждой
группе столбцов собственную таблицу. Тогда одинаковый номер
в двух группах уже не обязан означать одинаковую четвёрку весов.

Например, одна позиция может часто содержать небольшие
положительные веса, а другая --- сочетания другого масштаба
и знака. Независимые таблицы могут точнее описать
обе ситуации. Это мотивация, а не утверждение,
что реальные столбцы всегда имеют именно такой простой вид.

\question{Чем оплачивается дополнительная свобода?}
Индексов остаётся столько же, по одному на фрагмент,
но самих кодбуков становится намного больше.
Поэтому DPQ не заканчивается простым добавлением таблиц.
Нужны отдельное сжатие кодбуков и их дальнейшая калибровка.
Иначе небольшая номинальная битность индексов
будет скрывать большой реальный расход памяти.

Весь метод состоит из связанных частей: независимые
подпространства, общий пул близких кодовых векторов
внутри матрицы, перевыбор индексов по входным активациям,
обучение значений по задаче диффузии и поддержание
используемости центров. Воспроизведение только одной части
не является автоматически воспроизведением полного DPQ.

\slide{20}{Подпространства: что именно разделяется}
Подпространство здесь можно понимать как фиксированный
набор координат входного вектора. Для матрицы слоя
это одна и та же группа соседних столбцов во всех строках.
Это не группа изображений, не набор времён
и не отдельный трансформерный блок.

Пусть матрица имеет 1152 входных столбца, а длина фрагмента
равна четырём. Получается $1152/4=288$ групп.
Первая группа использует столбцы 1--4 всех строк,
вторая --- 5--8, последняя --- последние четыре столбца.
Для каждой из 288 групп строится собственный кодбук
из 256 векторов длины четыре.

Если первая часть строки выбирает $(1,2,3,4)$,
а вторая --- $(5,6,7,8)$, восстановленная строка
начинается как $(1,2,3,4,5,6,7,8)$.
Это конкатенация: векторы ставятся рядом.
Они не складываются и не перемножаются по координатам.
Название product относится к сочетаниям независимых выборов.

Для всех пяти матриц 28 блоков при $d=4$ получается
64\,512 логических кодбуков. Расчёт учитывает, что
у второй матрицы MLP 4608 входных столбцов,
а у остальных четырёх --- по 1152.
На блок: $(4\cdot1152+4608)/4=2304$;
на 28 блоков: $2304\cdot28=64\,512$.

\question{Чем полезно число $K^G$?}
Оно показывает количество возможных сочетаний выбора
по $G$ частям строки, не требуя хранить все строки.
Но само по себе это не доказательство превосходства
над общим кодбуком коротких фрагментов: у него тоже
много комбинаций индексов. Дополнительная свобода DPQ
связана с независимыми значениями таблиц разных групп.
Позднее часть этой свободы ограничивается объединением центров в пуле.

\slide{21}{Почему много кодбуков не бесплатно}
Число бит на индексы вычисляется так же, как прежде:
один номер из 256 требует восемь бит, а номер описывает
$d$ весов. Поэтому получается $8/d$ бита на вес.
При $d=2$ это четыре бита, при $d=4$ --- два,
при $d=8$ --- один. При $d=3$ получается $8/3$,
а не ровно три, хотя профиль в статье называется трёхбитным.

Однако у каждой группы собственные значения таблицы.
Для матрицы с $I$ входами, $O$ выходами и кодбуками FP16
число бит таблиц равно $(I/d)\cdot Kd\cdot16=IK\cdot16$.
Если разделить на число весов $OI$, добавка составит
$16K/O$ бита на вес. Она не исчезает из-за того,
что сами индексы короткие.

У матрицы $1152\times1152$, $K=256$, $d=4$
индексы требуют два бита на вес, а несжатые таблицы ---
ещё $16\cdot256/1152\approx3.56$.
Итого около 5.56 бита на вес до дополнительной компрессии,
без прочих частей модели. Это расчётная иллюстрация,
не опубликованный итоговый размер DPQ.

\question{Почему размер кодбуков почти не зависит от длины фрагмента?}
При уменьшении длины появляется больше групп,
но каждая запись становится короче. В этой конкретной
схеме с фиксированным $K$ два множителя сокращаются.
Меняется расход индексов, но несжатая общая таблица значений
не становится автоматически меньше.

\question{Какой вывод?}
Указывать только номинальные один-два бита недостаточно.
Нужно считать индексы, все кодовые векторы, дополнительные
ссылки и несжатые параметры. Для DPQ сжатие самих кодбуков
является необходимой частью метода, а не косметической оптимизацией файла.

\slide{22}{DPQ: близкие векторы хранятся один раз}
После независимого построения кодбуков разные группы
одной матрицы могут иметь похожие центры.
Возникает возможность не хранить одинаковые или почти
одинаковые наборы чисел много раз. Создаётся общий пул
векторов внутри одной матрицы и таблица ссылок на него.

На слайде приведён условный пример. Центр с номером 7
в первой группе равен $(0.10,0.20,0.30,0.40)$,
а центр с номером 19 во второй --- $(0.10,0.20,0.30,0.41)$.
Расстояние равно 0.01. При пороге 0.05 они достаточно
близки для объединения. В пуле остаётся одна запись,
а обе позиции кодбуков указывают на её номер 5.
Второй вектор теперь восстанавливается приближённо.

\question{Как восстанавливается конкретный фрагмент весов?}
Сначала его восьмибитный индекс выбирает запись
логического кодбука: например, 19 во второй группе.
Таблица проекции переводит пару «группа 2, запись 19»
в номер 5 общего пула. По номеру 5 читаются четыре
числа FP16. Итого два обращения: к таблице ссылок
и к значениям пула. Проекция здесь --- целочисленная
таблица, не обучаемая матрица линейного преобразования.

Часто используемые центры имеют приоритет при формировании
пула. Неиспользуемые или близкие записи не должны
получать такую же независимую память без необходимости.
У каждой матрицы свой пул: вся DiT не делит одну
глобальную таблицу со всеми слоями.

\question{Получается бесплатная компрессия?}
Нет. За объединение похожих, но разных векторов платится
ошибкой, а сама таблица ссылок тоже занимает память.
В рассматриваемом представлении ссылки 16-битные.
Нужно учитывать стоимость всей таблицы и пула,
а не демонстрировать только две удачно объединённые записи.
Детали формирования пула не полностью заданы опубликованным кодом;
выборы независимой реконструкции отдельно документированы.

\slide{23}{Калибровка начинается с настоящих изображений}
В DPQ источником данных служат реальные изображения
ImageNet с классами. Они переводятся VAE в чистые
латенты $z_0$. Затем выбирается время $t$ и создаётся
известный случайный массив $\varepsilon$.
Зашумлённый вход получается сразу по формуле
$z_t=\sqrt{\bar\alpha_t}z_0+\sqrt{1-\bar\alpha_t}\varepsilon$.

Первая часть сохраняет информацию исходного изображения,
вторая добавляет шум. Коэффициенты зависят от выбранного
времени. При слабом зашумлении заметнее изображение,
при сильном вход ближе к случайному шуму.
Шум известен, потому что он был создан при построении примера,
а не угадан из уже готового входа.

Квантованная DiT обрабатывает $z_t,t,y$.
По её входным признакам можно выбирать индексы,
сохраняющие локальные вычисления. Её окончательное
предсказание можно сравнивать с тем самым известным шумом
и учить значения кодбуков решать исходную задачу диффузии.

\question{Нужно ли предварительно выполнить сотни переходов?}
Нет. Один обучающий вход для произвольного времени
строится прямой формулой. На следующем примере выбирается
другое время и новый шум. Поэтому обучение охватывает
разные уровни зашумления без разворачивания всей обратной
генерации в каждом обновлении.

\question{Сеть обучается заново с нуля?}
Нет. Компактное представление получено из уже обученных
весов. Калибровка подстраивает его на исходной задаче.
Эталон M-шага --- известный шум, не предсказание
полноточного учителя и не готовая картинка.
Сама исходная матрица весов всё ещё нужна как
локальный ориентир при E-шаге. Это сочетание
послойной реконструкции и обучения на задаче,
а не обычная дистилляция ответов всей исходной сети.

\slide{24}{E-шаг: дискретно перевыбрать назначения}
На E-шаге значения центров фиксируются.
Для каждого фрагмента исходных весов нужно решить,
какой кодовый вектор текущей группы лучше использовать.
Но сравнение выполняется не только по расстоянию чисел
весов, а по ошибке их вклада в выход линейного слоя.

Пусть $X_j$ --- реальные входные признаки,
относящиеся к столбцам данной группы. Для кандидата
считается разница между исходным фрагментом
и центром $\delta=w-C_j[k]$. Ошибка равна
$\|X_j\delta\|^2$. Можно перебрать все 256
вариантов и сохранить номер с минимальной ошибкой.
Никакой производной через целое число для этого не требуется.

\question{Почему ближайший по весам центр может проиграть?}
Условный двумерный пример: первая координата входа
имеет средний квадрат 100, вторая --- 1,
а перекрёстная корреляция равна нулю.
У кандидата A ошибка весов $(0.1,0)$,
у B --- $(0,0.3)$. По весам A лучше:
$0.01<0.09$. Но ошибка выхода A равна
$100\cdot0.01=1$, B --- $1\cdot0.09=0.09$.
На этих входах выгоднее B. Это иллюстрация,
не измерение реального слоя DiT.

Матрица $H_j=X_j^\top X_j$ позволяет считать
ту же ошибку как $\delta^\top H_j\delta$.
В реальном случае она учитывает и корреляции признаков,
а не только отдельные масштабы.

\question{Это полная оптимизация выхода слоя?}
Нет. Группы рассматриваются независимо.
Перекрёстные члены между разными подпространствами
не оптимизируются совместно. Тем более это не
перебор индексов по итоговому FID или DDPM-loss всей сети.
Поэтому E-шаг может улучшать свою локальную цель,
но не даёт отдельной гарантии улучшения всей генерации.

\slide{25}{M-шаг: обновить кодбуки по DDPM-loss}
После выбора индексов они временно фиксируются.
Теперь восстановленные матрицы участвуют в полном
проходе квантованной DiT. Итоговое предсказание
шума сравнивается с известным случайным массивом,
который был добавлен к чистому латенту.

Получается обычная шумовая цель диффузии:
средний квадрат $\|\varepsilon-\varepsilon_{\widehat\theta}(z_t,t,y)\|^2$.
Градиент ошибки проходит через выходной модуль,
все блоки и восстановленные веса к значениям
использованных кодовых векторов. Если один вектор
использован многократно, его градиенты суммируются.
После сжатия таблиц обучаемыми значениями выступают
векторы общего пула соответствующей матрицы.

Смена чисел в кодбуках --- непрерывная оптимизация,
поэтому здесь подходит обычный AdamW. В статье
указан шаг $10^{-4}$ и по умолчанию одна эпоха
калибровки DiT: один проход по обучающим изображениям.
Для исходной DiT с обучаемой дисперсией стандартная
процедура также содержит вариационный компонент
для этой ветки. Формула на слайде выделяет
понятную основную цель по шуму, а не утверждает,
что дисперсия в реализации полностью игнорируется.

\question{Почему это не дистилляция?}
Для цели M-шага не нужен ответ исходной полной сети.
Метка --- известный добавленный шум. Это повторная
настройка компактного представления на задаче диффузии,
а не обучение повторять предсказание учителя.

\question{Что означает EM и есть ли гарантия улучшения?}
Здесь этим названием обозначено чередование
дискретных назначений и непрерывных значений.
E-шаг использует локальную ошибку подпространства,
M-шаг --- глобальную задачу сети. Это не точное
классическое EM с доказанной монотонностью одной
вероятностной цели. Улучшение каждого внутреннего
показателя не гарантирует уменьшения FID.
Пользу полного цикла необходимо измерять отдельно.

\slide{26}{Что делать с неиспользуемыми центрами}
Большой кодбук полезен только тогда, когда его
векторы действительно используются. Если ни один
фрагмент весов не ссылается на центр, через него
не проходит полезный градиент от восстановленных весов.
При фиксированных назначениях он не может сам
вернуться в работу только за счёт обновления
значений остальных центров.

В DPQ предусмотрено поддержание активных центров.
Для них ведётся возраст с последнего использования.
При новом назначении возраст обнуляется.
Если достаточно большая часть центров давно
не использовалась, запускается переинициализация.
На слайде приведены порог возраста 256
итераций и доля неактивных центров больше 3\%.
Это итерации процедуры калибровки,
не номера диффузионного времени изображения.

Для переинициализации нужен небольшой резервуар
примеров исходных векторов весов. Его не следует
путать с набором картинок ImageNet или с
кешем активаций. Живые центры объединяются
с содержимым резервуара, а новые положения
центров выбираются с помощью k-means++.
Затем настройка назначений и значений продолжается.

\question{Зачем нужны реальные примеры в резервуаре?}
Если новые центры выбирать только около оставшихся
активных, можно снова покрыть уже освоенную
область и не вернуть потерянные варианты весов.
Резервуар даёт дополнительные направления для
восстановления полезной ёмкости представления.

\question{Переинициализация полностью безвредна?}
Нет. Меняется представление, нужны согласованные
назначения и состояние оптимизатора.
Некоторые технические детали не однозначно
заданы авторским кодом и описанием статьи.
В независимой реконструкции они явно фиксируются.
Поддержание использования не является отдельным
набором кодбуков на каждое время и не увеличивает
автоматически итоговую битность.

\slide{27}{DPQ: опубликованные результаты}
Здесь приведена таблица 3 DPQ: DiT-XL/2,
ImageNet $256\times256$, DDPM250, CFG 1.5,
10\,000 изображений. Базовая FP16-модель
занимает 1349.4 МБ и имеет FID 5.41.
Это именно контроль данной таблицы,
а не значение 5.33 из другого авторского сравнения.

При двухбитном номинальном профиле авторский
VQ даёт FID 12.34, DPQ --- 6.84.
Размеры при этом различаются: 171.7
и 211.3 МБ. Значительная часть качества
возвращается, но дополнительное место
нельзя скрывать за одинаковой меткой «2 бита».

При одном бите разрыв ещё заметнее:
VQ имеет FID 91.71 и IS 13.51,
DPQ --- FID 14.03 и IS 110.23.
Размеры составляют 89.1 и 103.9 МБ.
Это показывает, что низкая номинальная битность
не обязательно полностью разрушает генерацию,
если представление и калибровка достаточно сильны.

\question{Получается ровно восьмикратное сжатие при двух битах?}
Нет. Реальная экономия относительно FP16
равна $1349.4/211.3\approx6.39$ раза.
При одном бите экономия составляет около 13 раз,
а не 16. Разницу создают кодбуки, ссылки
и несжатые части. Поэтому и показаны размеры,
а не только номинальные профили.

\question{Строка VQ является полной репродукцией VQ4DiT?}
Нет. Это вариант VQ, реализованный авторами DPQ.
Его нельзя без оговорки переименовать в
проверенный полный VQ4DiT. Для строгого сравнения
нужно понимать источник каждого базового метода,
условия его калибровки и реальную память.
Также эта таблица содержит опубликованные,
а не текущие собственные кластерные результаты.

\slide{28}{Два разных ответа на одну проблему}
У обеих статей общая цель: сохранить генерацию
при очень компактном представлении весов.
Но место основного изменения различается.
VQ4DiT оставляет один кодбук на матрицу
и делает назначение обучаемым через смесь кандидатов.
DPQ даёт группам столбцов отдельные таблицы,
затем ограничивает их память через общий пул.

В VQ4DiT источник калибровочных входов ---
промежуточные состояния генерации исходной DiT.
Эталон --- выходы её 28 блоков.
Изменяются значения кодбуков и непрерывные
предпочтения заранее выбранных кандидатов,
после чего остаётся один номер на фрагмент.

В DPQ источник --- реальные изображения
и прямое добавление известного шума.
E-шаг перебирает индексы по локальной
ошибке подпространства на реальных входах слоя.
M-шаг меняет значения по задаче
предсказания шума всей сетью.
Дополнительно поддерживается использование центров.

\question{Чем это отличается от текущего переноса VQDM?}
В VQDM фрагмент представлен суммой нескольких
кодовых векторов, а назначения уточняются
ограниченным совместным поиском. При исходной
подготовке используется локальная ошибка
слоя с учётом активаций. Это не тот же
набор кандидатов и не та же цель,
что в VQ4DiT или M-шаге DPQ.
Заимствование их обучающих сигналов остаётся
отдельной адаптацией, не репродукцией статьи.

\question{Какой вывод важнее рейтинга чисел?}
Недостаточно только приблизить веса по расстоянию.
Важны свобода назначения, распределение входов
и уровень обучающей цели. Однако сильные
опубликованные числа нужно подтверждать
собственным воспроизведением и согласованными
условиями: семплером, числом переходов,
VAE, фактическим размером и метриками.
Оба метода используют общие по времени
таблицы, а не отдельную модель для каждого шага.

\end{document}
